% !TEX root = ./main.tex

Adaptability can be intuitively understood as the property of a system
respond appropriately to changes in its execution environment.
%
Often, adaptation impacts on the behaviour of the system without
necessarily affecting its architecture.
%
For instance, fault-tolerance can be tackled by replicating components
without affecting the overall architecture.
%
If we embrace the view that for cloud applications this is not
necessarily the case, raises the problem of devising adaptation
mechanisms of architectures.
%
Arguably, the problem boils down to adopting suitable programming
features allowing developers to implement countermeasures against changes
in the execution context.
%
Let us introduce a running example that we use here to illustrate this
point and we will appeal to in order to demonstrate our approach throughout the
paper.
%
Our example elaborates on the adaptable
TeaStore benchmark\footnote{%
  These scenarios are based on the TeaStore use
  case~\cite{kistowski:ieee-mascots18} originally proposed to
  benchmark micro-services.} proposed
in~\cite{bliudze:arXiv-2412.16060} to evaluate adaptation mechanisms.
%
More precisely, this benchmark is designed to treat adaptability as a
first-class design concern in service-oriented architectures.
%
In particular, the benchmark classifies services as mandatory and
optional and as external or local, envisages multiple versions of
components (featuring e.g., different functionalities) in order to
specify and evaluate reconfiguration policies across diverse
operational scenarios.

\begin{description}
\item[Our running example] %\paragraph{Our running example.}
  %
  involves a web user-interface, an image provider service, and a
  locally hosted database service.
  % 
  The web user-interface, say \srv{WebUI}, is the entry point to the application
  for users; in particular it displays images of products generated by
  an image provider, say \srv{Image provider}.
  % 
  According to the specification in~\cite{kistowski:ieee-mascots18},
  \begin{itemize}
  \item image providers manipulate images (e.g., applying filters,
	 resizing, etc.);
  \item \srv{WebUI} connects to an image provider service in order to
	 generate webpages to display products to users.
  \end{itemize}
  A local storage service, say \srv{Local DB}, manages a database of images that
  \srv{WebUI} can resort to when e.g., the quality of images elaborated by
  \srv{Image provider} is sub-optimal.
  %
  The application can resort to locally stored default images if
  necessary.
\end{description}

Back to the application-dependent aspects of architectural adaptation,
let us consider a few possible ways to realize the reconfiguration
required when the TeaStore application needs to resort to the use of a
default image.
%
A first possibility is that service \textsf{WebUI} might trigger the search
for another image provider until the connection with the original
provider becomes acceptable.
% 
Alternatively, \textsf{WebUI} may resort to querying the service \textsf{Local DB}.
%
Or else, if the service \textsf{Local DB} performance is compromised by CPU overload, the \textsf{WebUI} might simply access a predefined file in the local drive, so avoiding further interactions with other services.
%
Obviously, which is the most reasonable options depends on the
application.

It is also worth remarking that, depending on the choice of how to
realize the adaptation, also the conditions triggering the
reconfiguration of the components in charge of the countermeasures to
adopt are application-dependent.
%
For instance, the second alternative could depend on the
responsiveness of \textsf{Local DB} when the adaptation is required.

\medskip

Although the actual tasks to execute when a reconfiguration is
required are application-dependent, there are activities that are
common to all cases.
%
In particular, the detection of anomalous situations that require
adaptation and the decision of which countermeasures to take can be
cast in the pattern presented in this paper.
%
Indeed, we explore the possibility of reusing known
mechanisms borrowed from the \emph{service-oriented computing} (SOC)
paradigm.\footnote{%
  Although the terminology has evolved, SOC is still widely used and,
  we believe, still applies to variants such as cloud computing, fog
  and edge computing, or the many forms of distributed computing
  associated with what is known as the \emph{Internet of (Every)Things}.
}
%
Our exploration hinges on a platform, dubbed \SEArch (after
{\ECFAugie S}ervice {\ECFAugie E}xecution {\ECFAugie
  Arch}itecture)~\cite{lopezpombo:coordination24}, for the discovery
and binding of distributed services featuring (non-)functional and
behavioural compliance of services.
%
We showcase how our platform can support architectural reconfigurations
dictated by  changes in the underlying
\emph{execution context}.
%
More precisely, in this paper we disregard adaptation that dependents on
application-level conditions and focus on reconfigurations due to
variations at lower levels such as network speed, CPU load, memory
consumption, etc.
%
As we will see, on the one hand this yields a simplified programming
pattern for adaptation and, on the other hand, is a limitation of our
approach.

\paragraph{Structure of the paper.}
\cref{sec:search} illustrates the key features of our platform
necessary to appreciate our programming pattern; more precisely, we
survey the behavioural contracts, and the searching and binding
mechanism featured by our platform.

\cref{sec:progpat} introduces our programming pattern and how it is cast
in the architecture of our platform.
%
We give the architecture of our adaptation framework in \cref{sec:arch};
the elements of the programming pattern are key to the design of the control loop
that decides how to reconfigure an application at run-time.
%
This decision procedure is detailed in \cref{sec:control-loop}
together with the specification of the adaptation contracts of our running example
which concern control data of the execution environment.

\cref{sec:teapat} describes the adaptation of our running example.
%
We give the behavioural contracts in terms of CFSMs for some scenarios
borrowed from the adaptable TeaStore
benchmark~\cite{bliudze:arXiv-2412.16060}.
%
Here, we also spell out how the adaptation process is triggered by the
control loop given in \cref{sec:control-loop}.

Finally, we draw some
conclusions and discuss further work in~\cref{sec:conc}.


% \begin{tikzpicture}[transform shape]
%   \node (broker) [draw, rectangle]{
% 	 \begin{minipage}{.2\linewidth}
% 		\textbf{Service Broker}
% 		\\
% 		\begin{tikzpicture}[transform shape]
% 		  \node (v) [draw, rectangle] {Pr $\vdash$ R};
% 		  \node (repo) [draw, rectangle, right = of v]{\textbf{Service Repository} \twemoji{computer disk}};
% 		\end{tikzpicture}
% 	 \end{minipage}
%   };
% \end{tikzpicture}


%%% Local Variables:
%%% mode: LaTeX
%%% TeX-master: "main"
%%% End:
